\documentclass[10pt,a4paper]{amsart}
\usepackage{amsmath,amssymb,amsthm,graphicx}
% Retained minimal preamble for the Stacks Project chapter sites-cohomology.tex (Cohomology on Sites).
% Verbatim excerpt of the chapter preamble: Xy-pic setup (native preamble lines 14-18) and the native
% theorem declarations with the mathematical macros (native preamble lines 154-217). Copied mechanically;
% nothing added, removed or redefined. The wrapper class and the amsmath/amssymb/amsthm/graphicx base
% packages are supplied by the assembler, so the class/externaldocument part of the native preamble is omitted.
% For commutative diagrams we use Xy-pic
\usepackage[all]{xy}

% We use 2cell for 2-commutative diagrams.
\xyoption{2cell}
% Theorem environments.
%
\theoremstyle{plain}
\newtheorem{theorem}[subsection]{Theorem}
\newtheorem{proposition}[subsection]{Proposition}
\newtheorem{lemma}[subsection]{Lemma}

\theoremstyle{definition}
\newtheorem{definition}[subsection]{Definition}
\newtheorem{example}[subsection]{Example}
\newtheorem{exercise}[subsection]{Exercise}
\newtheorem{situation}[subsection]{Situation}

\theoremstyle{remark}
\newtheorem{remark}[subsection]{Remark}
\newtheorem{remarks}[subsection]{Remarks}

\numberwithin{equation}{subsection}

% Macros
%
\def\lim{\mathop{\mathrm{lim}}\nolimits}
\def\colim{\mathop{\mathrm{colim}}\nolimits}
\def\Spec{\mathop{\mathrm{Spec}}}
\def\Hom{\mathop{\mathrm{Hom}}\nolimits}
\def\Ext{\mathop{\mathrm{Ext}}\nolimits}
\def\SheafHom{\mathop{\mathcal{H}\!\mathit{om}}\nolimits}
\def\SheafExt{\mathop{\mathcal{E}\!\mathit{xt}}\nolimits}
\def\Sch{\mathit{Sch}}
\def\Mor{\mathop{\mathrm{Mor}}\nolimits}
\def\Ob{\mathop{\mathrm{Ob}}\nolimits}
\def\Sh{\mathop{\mathit{Sh}}\nolimits}
\def\NL{\mathop{N\!L}\nolimits}
\def\CH{\mathop{\mathrm{CH}}\nolimits}
\def\proetale{{pro\text{-}\acute{e}tale}}
\def\etale{{\acute{e}tale}}
\def\QCoh{\mathit{QCoh}}
\def\Ker{\mathop{\mathrm{Ker}}}
\def\Im{\mathop{\mathrm{Im}}}
\def\Coker{\mathop{\mathrm{Coker}}}
\def\Coim{\mathop{\mathrm{Coim}}}

% Boxtimes
%
\DeclareMathSymbol{\boxtimes}{\mathbin}{AMSa}{"02}

%
% Macros for moduli stacks/spaces
%
\def\QCohstack{\mathcal{QC}\!\mathit{oh}}
\def\Cohstack{\mathcal{C}\!\mathit{oh}}
\def\Spacesstack{\mathcal{S}\!\mathit{paces}}
\def\Quotfunctor{\mathrm{Quot}}
\def\Hilbfunctor{\mathrm{Hilb}}
\def\Curvesstack{\mathcal{C}\!\mathit{urves}}
\def\Polarizedstack{\mathcal{P}\!\mathit{olarized}}
\def\Complexesstack{\mathcal{C}\!\mathit{omplexes}}
% \Pic is the operator that assigns to X its picard group, usage \Pic(X)
% \Picardstack_{X/B} denotes the Picard stack of X over B
% \Picardfunctor_{X/B} denotes the Picard functor of X over B
\def\Pic{\mathop{\mathrm{Pic}}\nolimits}
\def\Picardstack{\mathcal{P}\!\mathit{ic}}
\def\Picardfunctor{\mathrm{Pic}}
\def\Deformationcategory{\mathcal{D}\!\mathit{ef}}
\usepackage{cleveref}
\crefname{originallemma}{original lemma}{original lemmas}
\Crefname{originallemma}{Original lemma}{Original lemmas}
\makeatletter
\newlabel{lemma-build-from}{{0GYX}{source}{}{sourcepointer-1}{}}
\newlabel{lemma-build-from@cref}{{[originallemma][0][]0GYX}{[1][0][]source}{}{}{}}
\makeatother
\crefname{originallemma}{original lemma}{original lemmas}
\Crefname{originallemma}{Original lemma}{Original lemmas}
\makeatletter
\newlabel{lemma-qis-small}{{0GYY}{source}{}{sourcepointer-2}{}}
\newlabel{lemma-qis-small@cref}{{[originallemma][0][]0GYY}{[1][0][]source}{}{}{}}
\makeatother
\crefname{originallemma}{original lemma}{original lemmas}
\Crefname{originallemma}{Original lemma}{Original lemmas}
\makeatletter
\newlabel{more-algebra-lemma-enlarge}{{0GYI}{source}{}{sourcepointer-3}{}}
\newlabel{more-algebra-lemma-enlarge@cref}{{[originallemma][0][]0GYI}{[1][0][]source}{}{}{}}
\makeatother
\title{In the quasi-coherent situation of a ringed site there is a cardinal kappa bounding the size of a nonzero map to any nonzero object and of a factorisation through a direct sum}
\author{The Stacks Project Authors}
\date{stacks-project:sites-cohomology.tex:local-snapshot-20261011; GNU FDL 1.2 or later}
\begin{document}
\maketitle
\noindent\textit{Editorial note.} All delivered mathematical text is the author's own, copied from the pinned Stacks Project chapter file sites-cohomology.tex; no statement or proof step was rewritten, reordered or removed.\par\medskip
\noindent\textit{Editorial note.} The counted proof is the author's complete argument: the proof is the author's complete argument: the three-cardinal definition of kappa, the reduction of both assertions to the claim, and the induction proving the claim with the choice of the subcomplexes N\_f, M\_f and the application of lemma-build-from are all written out, ending with the colimit argument that H represents an object of QC(O).\par\medskip
\noindent\textit{Editorial note.} The label lemma-cartesian-colimit-kappa names the counted statement, and the proof environment immediately following it in the source is delivered verbatim as its proof.\par\medskip
\noindent\textit{Editorial note.} The statement and proof are copied from the pinned chapter file; no step was reordered, merged or removed.\par\medskip
\noindent\textit{Editorial note.} The preamble is the chapter's own retained minimal preamble (sources/stacks-sites-cohomology/retained-preamble.tex), a verbatim excerpt of the native preamble: the Xy-pic setup and the native theorem declarations with the mathematical macros.\par\medskip
\noindent\textit{Allowed standard prerequisites (editorial).}\begin{itemize}
\item The category QC(O) of quasi-coherent objects of D(O). Same chapter, the full subcategory QC(O) of D(O) defined by the condition that the derived tensor comparison maps are isomorphisms and its presentation by complexes. Supplies the ambient category in which both properties are stated and in which the complexes of size <= kappa are chosen.
\item Size |K| of an object of QC(O). Same chapter, the cardinality |K| of an object of QC(O) defined as the minimum size of a complex of O-modules representing it. Supplies the quantity bounded by kappa in both assertions of the statement and in the inductive construction.
\item The quasi-compact, quasi-separated local situation of the section. Same chapter, the standing situation above the lemma (the site C with its structural sheaf O and the auxiliary objects j\_\{U!\}O\_U). Supplies the coproducts appearing in the definition of kappa and the objects used in part (1).
\end{itemize}\medskip
\ifdefined\hypertarget\hypertarget{sourcepointer-1}{}\fi
\noindent\textit{Source reference.} Original lemma 0GYX: sources/stacks-sites-cohomology/pinned-original-source.tex, source lines 11217--11217. Prerequisite pointer invoked at proof line(s) 11434; it names a standard fact of the same chapter that the counted argument uses, and every step of the counted proof built on it is retained verbatim in the delivered body.\par\medskip
\ifdefined\hypertarget\hypertarget{sourcepointer-2}{}\fi
\noindent\textit{Source reference.} Original lemma 0GYY: sources/stacks-sites-cohomology/pinned-original-source.tex, source lines 11238--11238. Prerequisite pointer invoked at proof line(s) 11318; it names a standard fact of the same chapter that the counted argument uses, and every step of the counted proof built on it is retained verbatim in the delivered body.\par\medskip
\ifdefined\hypertarget\hypertarget{sourcepointer-3}{}\fi
\noindent\textit{Source reference.} Original lemma 0GYI: sources/stacks-chow/tags.txt, source lines 21143--21143. Prerequisite pointer invoked at proof line(s) 11316; it names a standard fact of a sibling chapter that the counted argument uses, and every step of the counted proof built on it is retained verbatim in the delivered body.\par\medskip
\begin{lemma}
\label{lemma-cartesian-colimit-kappa}
In the situation above, there exists a cardinal $\kappa$ with the following
properties:
\begin{enumerate}
\item for every nonzero object $K$ of $\mathit{QC}(\mathcal{O})$ there exists a
nonzero morphism $E \to K$ of $\mathit{QC}(\mathcal{O})$ such that
$|E| \leq \kappa$,
\item for every morphism $\alpha : E \to \bigoplus_n K_n$ of
$\mathit{QC}(\mathcal{O})$ such that $|E| \leq \kappa$, there
exist morphisms $E_n \to K_n$
in $\mathit{QC}(\mathcal{O})$ with $|E_n| \leq \kappa$
such that $\alpha$ factors through $\bigoplus E_n \to \bigoplus K_n$.
\end{enumerate}
\end{lemma}

\begin{proof}
Let $\kappa$ be an upper bound for the following set of cardinals:
\begin{enumerate}
\item $|\coprod_V j_{U!}\mathcal{O}_U(V)|$ for all $U \in \Ob(\mathcal{C})$,
\item the cardinals $\kappa(\mathcal{O}(V) \to \mathcal{O}(U))$
found in More on Algebra, Lemma \ref{more-algebra-lemma-enlarge}
for all morphisms $U \to V$ in $\mathcal{C}$,
\item the cardinal found in Lemma \ref{lemma-qis-small}.
\end{enumerate}
We claim that for any complex $\mathcal{F}^\bullet$ representing an object of
$\mathit{QC}(\mathcal{O})$ and any subcomplex
$\mathcal{S}^\bullet \subset \mathcal{F}^\bullet$
with $|\mathcal{S}^\bullet| \leq \kappa$ there exists a subcomplex
$\mathcal{H}^\bullet$ of $\mathcal{F}^\bullet$ containing $\mathcal{S}^\bullet$
such that $\mathcal{H}^\bullet$ represents an object of
$\mathit{QC}(\mathcal{O})$
and such that $|\mathcal{H}^\bullet| \leq \kappa$.
In the next two paragraphs we show that the claim implies the lemma.

\medskip\noindent
As in (1) let $K$ be a nonzero object of $\mathit{QC}(\mathcal{O})$.
Say $K$ is represented by the complex
of $\mathcal{O}$-modules $\mathcal{F}^\bullet$. Then $H^i(\mathcal{F}^\bullet)$
is nonzero for some $i$. Hence there exists an object $U$ of $\mathcal{C}$
and a section $s \in \mathcal{F}^i(U)$ with $d(s) = 0$ which determines a
nonzero section of $H^i(\mathcal{F}^\bullet)$ over $U$. Then the image
of $s : j_{U!}\mathcal{O}_U[-i] \to \mathcal{F}^\bullet$ is a subcomplex
$\mathcal{S}^\bullet \subset \mathcal{F}^\bullet$ with
$|\mathcal{S}^\bullet| \leq \kappa$. Applying the claim we get
$\mathcal{H}^\bullet \to \mathcal{F}^\bullet$ in
$\mathit{QC}(\mathcal{O})$ nonzero
with $|\mathcal{H}^\bullet| \leq \kappa$. Thus (1) holds.

\medskip\noindent
Let $\alpha : E \to \bigoplus K_n$ be as in (2).
Choose any complexes $\mathcal{K}_n^\bullet$ representing $K_n$.
Then $\bigoplus \mathcal{K}_n^\bullet$ represents $\bigoplus K_n$.
By the construction of the derived category we can represent
$E$ by a complex $\mathcal{E}^\bullet$ such that $\alpha$
is represented by a morphism
$a : \mathcal{E}^\bullet \to \bigoplus \mathcal{K}_n^\bullet$
of complexes. By Lemma \ref{lemma-qis-small} and our choice of
$\kappa$ above we may assume
$|\mathcal{E}^\bullet| \leq \kappa$. By the claim
we get subcomplexes $\mathcal{E}_n^\bullet \subset \mathcal{K}_n^\bullet$
representing objects $E_n$ of
$\mathit{QC}(\mathcal{O})$ with $|E_n| \leq \kappa$
containing the image of $a_n : \mathcal{E}^\bullet \to \mathcal{K}_n^\bullet$
as desired.

\medskip\noindent
Proof of the claim. Let $\mathcal{F}^\bullet$ be a complex representing
an object of $\mathit{QC}(\mathcal{O})$ and let
$\mathcal{S}^\bullet \subset \mathcal{F}^\bullet$ be a subcomplex of
size $\leq \kappa$. We are going to inductively construct subcomplexes
$$
\mathcal{S}^\bullet = \mathcal{S}_0^\bullet \subset
\mathcal{S}_1^\bullet \subset
\mathcal{S}_2^\bullet \subset \ldots \subset \mathcal{F}^\bullet
$$
of size $\leq \kappa$ such that for every morphism $f : U \to V$
of $\mathcal{C}$ and every $i \in \mathbf{Z}$
\begin{enumerate}
\item the kernel of the arrow
$H^i(\mathcal{S}_n^\bullet(V)
\otimes_{\mathcal{O}(V)}^\mathbf{L} \mathcal{O}(U)) \to
H^i(\mathcal{S}_n^\bullet(U))$
maps to zero in
$H^i(\mathcal{S}_{n + 1}^\bullet(V)
\otimes_{\mathcal{O}(V)}^\mathbf{L} \mathcal{O}(U))$,
\item the image of the arrow
$H^i(\mathcal{S}_n^\bullet(U)) \to H^i(\mathcal{S}_{n + 1}^\bullet(U))$
is contained in the image of
$H^i(\mathcal{S}_{n + 1}^\bullet(V)
\otimes_{\mathcal{O}(V)}^\mathbf{L} \mathcal{O}(U)) \to
H^i(\mathcal{S}_{n + 1}^\bullet(U))$,
\end{enumerate}
Once this is done we can set
$\mathcal{H}^\bullet = \bigcup \mathcal{S}_n^\bullet$.
Namely, since derived tensor product and taking cohomology of
complexes of modules over rings commute with filtered
colimits, the conditions (1) and (2) together will guarantee
that
$$
\mathcal{H}^\bullet(V) \otimes_{\mathcal{O}(V)}^\mathbf{L} \mathcal{O}(U)
\longrightarrow
\mathcal{H}^\bullet(U)
$$
is an isomorphism on cohomology in all degrees and hence an isomorphism
in $D(\mathcal{O}(U))$ for all $f : U \to V$ in $\mathcal{C}$.
Hence $\mathcal{H}^\bullet$ represents an object of
$\mathit{QC}(\mathcal{O})$ as desired.

\medskip\noindent
Construction of $\mathcal{S}_{n + 1}$ given $\mathcal{S}_n$.
For every morphism $f : U \to V$ of $\mathcal{C}$
we consider the commutative diagram
$$
\xymatrix{
\mathcal{S}_n^\bullet(V)
\ar[r] \ar[d] &
\mathcal{S}_n^\bullet(U) \ar[d] \\
\mathcal{F}^\bullet(V)
\ar[r] &
\mathcal{F}^\bullet(U)
}
$$
This is a diagram as in More on Algebra, Lemma \ref{more-algebra-lemma-enlarge}
for the ring map $\mathcal{O}(V) \to \mathcal{O}(U)$, i.e., the bottom row
induces an isomorphism
$$
\mathcal{F}^\bullet(V) \otimes_{\mathcal{O}(V)}^\mathbf{L} \mathcal{O}(U)
\longrightarrow
\mathcal{F}^\bullet(U)
$$
in $D(\mathcal{O}(U))$. Thus we may choose subcomplexes
$$
\mathcal{S}_n^\bullet(V) \subset M^\bullet_f \subset \mathcal{F}^\bullet(V)
\quad\text{and}\quad
\mathcal{S}_n^\bullet(U) \subset N^\bullet_f \subset \mathcal{F}^\bullet(U)
$$
as in More on Algebra, Lemma \ref{more-algebra-lemma-enlarge}
and in particular we see that $|N^i_f|, |M^i_f| \leq \kappa$. Next, we apply
Lemma \ref{lemma-build-from} using the subsets
$$
\mathcal{S}_n^i(U) \amalg \coprod\nolimits_{f : U \to V} N^i_f
\amalg \coprod\nolimits_{g : W \to U} M^i_g
\subset
\mathcal{F}^i(U)
$$
to find a subcomplex
$$
\mathcal{S}_n^\bullet \subset
\mathcal{S}_{n + 1}^\bullet \subset
\mathcal{F}^\bullet
$$
with containing those subsets and such that
$|\mathcal{S}_{n + 1}^\bullet| \leq \kappa$.
Conditions (1) and (2) hold because the corresponding statements
hold for 
$\mathcal{S}_n^\bullet(V) \subset M^\bullet_f$ and
$\mathcal{S}_n^\bullet(U) \subset N^\bullet_f$ by the construction in
More on Algebra, Lemma \ref{more-algebra-lemma-enlarge}.
Thus the proof is complete.
\end{proof}

\end{document}
